% GENERATED FILE. Edit canonical lessons, then run npm run generate:books.
% canonical-source-hash: fnv1a64:28f94ddd89258798
% canonical-lessons: KA-C202-duru-kodu, KA-C202-nondayisu, KA-C202-munduvarisu, KA-C202-samparkisu, KA-C202-carcisu

\chapter{To complain, To register, To continue, To contact, To discuss}
\label{ch:ka-to-complain-to-register-to-continue-to-contact-to-discuss}
\hlchaptermodality{202}

\begin{chapteropening}
\textbf{By the end of this chapter:} \emph{I can say to complain, to register, to continue, to contact and to discuss.}

Complete the last lesson of chapter 202: I can say to complain, to register, to continue, to contact and to discuss.
\end{chapteropening}

\section[\texorpdfstring{\kn{ದೂರು} \kn{ಕೊಡು} (dūru koḍu)}{ದೂರು ಕೊಡು (dūru koḍu)}]{\kn{ದೂರು} \kn{ಕೊಡು} (dūru koḍu) — to complain, to lodge a complaint}

\label{lesson:KA-C202-duru-kodu}



\begin{quote}
Before the new one: say the Kannada for to remember, then the Kannada for to print.
\end{quote}

\subsection*{You'll want to know: \kn{ದೂರು} \kn{ಕೊಡು}}
\textbf{\kn{ದೂರು} \kn{ಕೊಡು}} — \emph{dūru koḍu} — \textquotedblleft{}to complain, to lodge a complaint\textquotedblright{}.

\begin{grammarlens}[title={What you've built}]
One more word for this chapter.
\end{grammarlens}

\subsection*{Your turn}
\begin{itemize}
\raggedright
  \item \emph{Say it:} \emph{dūru koḍu}
  \item \emph{Say it:} \emph{dūru koḍu}, once more
  \item \emph{Say it:} say \emph{dūru koḍu}
  \item \emph{Recall:} say the Kannada for to remember, then the Kannada for to print, then say \emph{dūru koḍu} again
\end{itemize}

\begin{tcolorbox}[breakable,colback=teal!4,colframe=teal!35!black,title={Before you move on}]
What does \kn{ದೂರು} \kn{ಕೊಡು} mean? (\textquotedblleft{}To complain, to lodge a complaint\textquotedblright{}.) Say it once more.
\end{tcolorbox}
\section[\texorpdfstring{\kn{ನೋಂದಾಯಿಸು} (nōndāyisu)}{ನೋಂದಾಯಿಸು (nōndāyisu)}]{\kn{ನೋಂದಾಯಿಸು} (nōndāyisu) — to register}

\label{lesson:KA-C202-nondayisu}



\begin{quote}
Before the new one: say the Kannada for to print, then the Kannada for to complain.
\end{quote}

\subsection*{You'll want to know: \kn{ನೋಂದಾಯಿಸು}}
\textbf{\kn{ನೋಂದಾಯಿಸು}} — \emph{nōndāyisu} — \textquotedblleft{}to register\textquotedblright{}.

\begin{grammarlens}[title={What you've built}]
One more word for this chapter.
\end{grammarlens}

\subsection*{Your turn}
\begin{itemize}
\raggedright
  \item \emph{Say it:} \emph{nōndāyisu}
  \item \emph{Say it:} \emph{nōndāyisu}, once more
  \item \emph{Say it:} say \emph{nōndāyisu}
  \item \emph{Recall:} say the Kannada for to print, then the Kannada for to complain, then say \emph{nōndāyisu} again
\end{itemize}

\begin{tcolorbox}[breakable,colback=teal!4,colframe=teal!35!black,title={Before you move on}]
What does \kn{ನೋಂದಾಯಿಸು} mean? (\textquotedblleft{}To register\textquotedblright{}.) Say it once more.
\end{tcolorbox}
\section[\texorpdfstring{\kn{ಮುಂದುವರಿಸು} (munduvarisu)}{ಮುಂದುವರಿಸು (munduvarisu)}]{\kn{ಮುಂದುವರಿಸು} (munduvarisu) — to continue}

\label{lesson:KA-C202-munduvarisu}



\begin{quote}
Before the new one: say the Kannada for to complain, then the Kannada for to register.
\end{quote}

\subsection*{You'll want to know: \kn{ಮುಂದುವರಿಸು}}
\textbf{\kn{ಮುಂದುವರಿಸು}} — \emph{munduvarisu} — \textquotedblleft{}to continue\textquotedblright{}.

\begin{grammarlens}[title={What you've built}]
One more word for this chapter.
\end{grammarlens}

\subsection*{Your turn}
\begin{itemize}
\raggedright
  \item \emph{Say it:} \emph{munduvarisu}
  \item \emph{Say it:} \emph{munduvarisu}, once more
  \item \emph{Say it:} say \emph{munduvarisu}
  \item \emph{Recall:} say the Kannada for to complain, then the Kannada for to register, then say \emph{munduvarisu} again
\end{itemize}

\begin{tcolorbox}[breakable,colback=teal!4,colframe=teal!35!black,title={Before you move on}]
What does \kn{ಮುಂದುವರಿಸು} mean? (\textquotedblleft{}To continue\textquotedblright{}.) Say it once more.
\end{tcolorbox}
\section[\texorpdfstring{\kn{ಸಂಪರ್ಕಿಸು} (samparkisu)}{ಸಂಪರ್ಕಿಸು (samparkisu)}]{\kn{ಸಂಪರ್ಕಿಸು} (samparkisu) — to contact}

\label{lesson:KA-C202-samparkisu}



\begin{quote}
Before the new one: say the Kannada for to register, then the Kannada for to continue.
\end{quote}

\subsection*{You'll want to know: \kn{ಸಂಪರ್ಕಿಸು}}
\textbf{\kn{ಸಂಪರ್ಕಿಸು}} — \emph{samparkisu} — \textquotedblleft{}to contact\textquotedblright{}.

\begin{grammarlens}[title={What you've built}]
One more word for this chapter.
\end{grammarlens}

\subsection*{Your turn}
\begin{itemize}
\raggedright
  \item \emph{Say it:} \emph{samparkisu}
  \item \emph{Say it:} \emph{samparkisu}, once more
  \item \emph{Say it:} say \emph{samparkisu}
  \item \emph{Recall:} say the Kannada for to register, then the Kannada for to continue, then say \emph{samparkisu} again
\end{itemize}

\begin{tcolorbox}[breakable,colback=teal!4,colframe=teal!35!black,title={Before you move on}]
What does \kn{ಸಂಪರ್ಕಿಸು} mean? (\textquotedblleft{}To contact\textquotedblright{}.) Say it once more.
\end{tcolorbox}
\section[\texorpdfstring{\kn{ಚರ್ಚಿಸು} (carcisu)}{ಚರ್ಚಿಸು (carcisu)}]{\kn{ಚರ್ಚಿಸು} (carcisu) — to discuss}

\label{lesson:KA-C202-carcisu}



\begin{quote}
Before the new one: say the Kannada for to continue, then the Kannada for to contact.
\end{quote}

\subsection*{You'll want to know: \kn{ಚರ್ಚಿಸು}}
\textbf{\kn{ಚರ್ಚಿಸು}} — \emph{carcisu} — \textquotedblleft{}to discuss\textquotedblright{}.

\begin{grammarlens}[title={What you've built}]
One more word for this chapter.
\end{grammarlens}

\subsection*{Your turn}
\begin{itemize}
\raggedright
  \item \emph{Say it:} \emph{carcisu}
  \item \emph{Say it:} \emph{carcisu}, once more
  \item \emph{Say it:} say \emph{carcisu}
  \item \emph{Recall:} say the Kannada for to continue, then the Kannada for to contact, then say \emph{carcisu} again
\end{itemize}

\begin{tcolorbox}[breakable,colback=teal!4,colframe=teal!35!black,title={Before you move on}]
What does \kn{ಚರ್ಚಿಸು} mean? (\textquotedblleft{}To discuss\textquotedblright{}.) Say it once more.
\end{tcolorbox}
